% 由 工程文件\tools\gen_tables.py 生成，勿手改。
% 译注数据：工程文件\tools\notes_data.py

\expandafter\def\csname zhnlab@1.1\endcsname{1.1}
\expandafter\def\csname zhnbody@1.1\endcsname{$\YC(G,C)$ 是 $C$ 上平坦 $G_{\C}$-丛的模空间；等价地，是 $\pi_{1}(C)\to G_{\C}$ 的同态模去共轭所得的表示簇。它一般并不光滑：可约表示（像落在真子群中的那些）构成若干真子簇，下文谈镜像对称时实际上只在不可约部分上工作，那里它是复流形，复维数为 $(g-1)\dim G$。又，$\YC$ 只依赖 $C$ 的拓扑（基本群），这正是后文反复强调「对偶陈述只依赖 $C$ 的定向二维流形结构」的起点。}
\expandafter\def\csname zhnsrc@1.1\endcsname{[来源：B；nLab 词条「flat connection」（2026-11-15 访问）；Atiyah–Bott (1982)]}

\expandafter\def\csname zhnlab@1.2\endcsname{1.2}
\expandafter\def\csname zhnbody@1.2\endcsname{$\Omega$ 是 $\YC(G,C)$ 上的全纯辛形式。规范场论的写法是 $\Omega=\int_{C}\operatorname{tr}(\delta A\wedge\delta A)$，其中 $\delta A$ 是联络的变分（取值于 $\ad(E)$）；Atiyah 与 Bott 证明它非退化，并把 $\YC(G,C)$ 实现为无穷维辛商 $\mathcal{A}/\mathcal{G}_{\C}$ 的辛约化。这种「先在无穷维仿射空间上取商、再把辛形式下降」的套路，在第 5、6 节会以超凯勒商的形式再次出现。}
\expandafter\def\csname zhnsrc@1.2\endcsname{[来源：A+B；nLab 词条「symplectic quotient」（Last revised September 20, 2013；2026-11-15 访问）；Atiyah–Bott (1982)]}

\expandafter\def\csname zhnlab@1.3\endcsname{1.3}
\expandafter\def\csname zhnbody@1.3\endcsname{典范线丛平凡，即存在处处非零的全纯体积形式 $\Omega^{n}/n!$（$2n=\dim_{\C}\YC$）。它的作用是：复结构的选择只影响这个体积形式的整体相位，而 A-模型的计算（伪全纯曲线的计数）与这一相位无关，于是 A-模型确实只依赖 $\omega=\operatorname{Im}\Omega$。对典范丛非平凡的紧 Kähler 流形，A-模型会因 $c_{1}\neq0$ 而带上对复结构的依赖，原文这句断语便不再成立。}
\expandafter\def\csname zhnsrc@1.3\endcsname{[来源：B；nLab 词条「A-model」（Last revised October 23, 2026；2026-11-15 访问）；原文脚注 3]}

\expandafter\def\csname zhnlab@1.4\endcsname{1.4}
\expandafter\def\csname zhnbody@1.4\endcsname{A-模型对辛形式的整体标度不敏感：$\Omega\mapsto\lambda\Omega$ 使 $\omega\mapsto\lambda\omega$，而关联函数只依赖 $\omega$ 的复化上同调类，整体因子可被吸收进重标度。因此「用 $\operatorname{Re}\Omega$ 还是 $\operatorname{Im}\Omega$」并非无关紧要：$\operatorname{Re}\Omega_{I}=\omega_{J}$ 与 $\operatorname{Im}\Omega_{I}=\omega_{K}$ 是同一个超凯勒 $S^{2}$ 族中的两个不同点，取不同的实线性组合就得到不同的 A-模型——这正是下一段「量子几何 Langlands」的参数来源。}
\expandafter\def\csname zhnsrc@1.4\endcsname{[来源：B；nLab 词条「A-model」（Last revised October 23, 2026；2026-11-15 访问）；nLab 词条「mirror symmetry」（2026-11-15 访问）；原文脚注 3；Kapustin–Witten (2007) 第 5.2 节]}

\expandafter\def\csname zhnlab@1.5\endcsname{1.5}
\expandafter\def\csname zhnbody@1.5\endcsname{「量子几何 Langlands」指几何 Langlands 的一个单参数（更准确地说是 $S^{2}$ 值）推广：在 A-模型一侧把辛形式换成 $\operatorname{Im}(e^{i\vartheta}\Omega)$，$\vartheta$ 跑遍 $S^{2}$；B-模型一侧则相应地「形变」为一个不再是通常复流形的理论。之所以能这样做，是因为这种形变在广义复几何中有精确含义：广义复结构把复结构与辛结构放进同一个 $S^{2}$ 族里插值。Kapustin 与 Witten 说明，这个推广的参数正是四维 $\mathcal{N}=4$ 理论中 $\theta$-角与耦合常数的组合。}
\expandafter\def\csname zhnsrc@1.5\endcsname{[来源：B；nLab 词条「generalized complex geometry」（2026-11-15 访问）；Hitchin (2003)；Gualtieri (2003)；Kapustin–Witten (2007) 第 5.2、11.3 节]}

\expandafter\def\csname zhnlab@1.6\endcsname{1.6}
\expandafter\def\csname zhnbody@1.6\endcsname{${}^{L}G$ 是 $G$ 的 Langlands 对偶群：取 $G$ 的根数据，把根与余根互换（等价地把权格与余权格互换）所得的群。例：$\SU(N)$ 的对偶是 $\PSU(N)=\SU(N)/\mathbb{Z}_{N}$，环面的对偶是取对偶环面（角色群），$E_{8}$ 自对偶。物理上它就是 Montonen–Olive 电磁对偶中「磁」一侧的规范群。注意 ${}^{L}G$ 与 $G$ 一般不同构，它们的复化也不同（如 $\PSU(N)_{\C}=\mathrm{PGL}(N,\C)$）。}
\expandafter\def\csname zhnsrc@1.6\endcsname{[来源：B；nLab 词条「Langlands dual group」（2026-11-15 访问）；Montonen–Olive (1977)；Goddard–Nuyts–Olive, Nucl. Phys. B125 (1977) 1]}

\expandafter\def\csname zhnlab@2.1\endcsname{2.1}
\expandafter\def\csname zhnbody@2.1\endcsname{镜像对称是两台二维拓扑场论之间的等价：流形 $X$ 的 A-模型（辛几何侧）等价于另一个流形 $X'$ 的 B-模型（复几何侧）。它不是经典几何的等同，而是「量子」层面的等同，这正是原文称之为 quantum symmetry 的意思；同调镜对称把它实现为 $\mathcal{D}^{b}\mathrm{Fuk}(X)\simeq\mathcal{D}^{b}\mathrm{Coh}(X')$ 的导出等价。几何图像由 Strominger–Yau–Zaslow 给出：$X$ 与 $X'$ 是同一个底上互为对偶的特殊拉格朗日环面纤维化。}
\expandafter\def\csname zhnsrc@2.1\endcsname{[来源：B；nLab 词条「mirror symmetry」（2026-11-15 访问）；nLab 词条「A-model」（Last revised October 23, 2026；2026-11-15 访问）；nLab 词条「B-model」（Last revised February 12, 2025；2026-11-15 访问）；Strominger–Yau–Zaslow (1996)]}

\expandafter\def\csname zhnlab@2.2\endcsname{2.2}
\expandafter\def\csname zhnbody@2.2\endcsname{之所以猜想 $\YC(\LG,C)$ 的 B-模型镜像于 $\YC(G,C)$ 的 A-模型（而不是反过来），原因在于几何 Langlands 的惯例：对偶群一侧承载「磁」的、复几何的数据（局部系统、Hecke 特征层），原群一侧承载「电」的、辛几何的数据。Hausel–Thaddeus 正是以此猜想为出发点，用 Hitchin 纤维化给出这一镜像对称的第一个具体图像；Kapustin–Witten 则把它实现为四维 $\mathcal{N}=4$ 理论的 S-对偶。}
\expandafter\def\csname zhnsrc@2.2\endcsname{[来源：A+B；nLab 词条「Hitchin fibration」（2026-11-15 访问）；nLab 词条「geometric Langlands correspondence」（Last revised April 11, 2026；2026-11-15 访问）；Hausel–Thaddeus (2002)]}

\expandafter\def\csname zhnlab@2.3\endcsname{2.3}
\expandafter\def\csname zhnbody@2.3\endcsname{式 (2.1) 是三组方程。前两式合起来等价于复联络 $\mathcal{A}=A+i\phi$ 的平坦性：$F(\mathcal{A})=F-\phi\wedge\phi+iD\phi=0$，即 $F-\phi\wedge\phi=0$ 与 $D\phi=0$。第三式 $D\star\phi=0$ 不在平坦性里，它是额外的一个实方程（$\star$ 为 $C$ 上复结构决定的 Hodge 星算子）；正是它使解模空间成为超凯勒流形，而不只是复辛流形（只取前两式得到 $\YC(G,C)$，那里只有复辛结构）。这也是原文说「$\star$ 的作用是让最后一个方程可写」的确切含义。}
\expandafter\def\csname zhnsrc@2.3\endcsname{[来源：A+B；nLab 词条「Hitchin system」（2026-11-15 访问）；nLab 词条「Bogomolny equation」（2026-11-15 访问）；Hitchin (1987a)]}

\expandafter\def\csname zhnlab@2.4\endcsname{2.4}
\expandafter\def\csname zhnbody@2.4\endcsname{Higgs 丛是 (2.1) 的解在复结构 $I$ 下的全纯翻译：一个全纯 $G_{\C}$-丛 $E$ 连同 $\phi\in H^{0}(C,K\otimes\ad(E))$。名称来自规范场论中的 Higgs 机制：$\phi$ 扮演那个赋予质量的标量场，只是此处它取值于 $\ad(E)\otimes K$。Hitchin 与 Simpson 证明：$E$ 上存在满足 Hitchin 方程的联络，当且仅当 $(E,\phi)$ 是（斜率）稳定的，且解模去规范变换与稳定 Higgs 丛的模空间同构——这就是非阿贝尔 Hodge 对应。$\MH(G,C)$ 的复维数是 $2\dim_{\C}M(G,C)=2(g-1)\dim G$。}
\expandafter\def\csname zhnsrc@2.4\endcsname{[来源：A+B；nLab 词条「Higgs bundle」（Last revised August 12, 2026；2026-11-15 访问）；Hitchin (1987a)；Simpson (1990)]}

\expandafter\def\csname zhnlab@2.5\endcsname{2.5}
\expandafter\def\csname zhnbody@2.5\endcsname{超凯勒流形是带三个复结构 $I,J,K$（满足 $IJ=K$、$I^{2}=J^{2}=K^{2}=-1$ 的四元数关系）的流形，且相应的三个 Kähler 形式 $\omega_{I},\omega_{J},\omega_{K}$ 都闭。Hitchin 证明 $\MH$ 是超凯勒的，其构造正是第 6.1 节末尾所说的超凯勒商。关键推论：同一个光滑流形上同时住着三个复结构，于是 $(I,\Omega_{I})$ 给出 B-模型，$(J,\Omega_{J})$ 给出另一个 B-模型，而 $\omega_{K}=\operatorname{Im}\Omega_{I}$ 给出 A-模型——「同一空间既是 A 又是 B」之所以可能，全靠这个结构。}
\expandafter\def\csname zhnsrc@2.5\endcsname{[来源：A+B；nLab 词条「hyperkähler manifold」（Last revised July 14, 2021；2026-11-15 访问）；Hitchin (1987a)；Hitchin–Karlhede–Lindström–Roček (1987)]}

\expandafter\def\csname zhnlab@2.6\endcsname{2.6}
\expandafter\def\csname zhnbody@2.6\endcsname{在复结构 $J$ 下，$\MH(G,C)$ 就是 $\YC(G,C)$：因为 $J$ 是「把联络与 Higgs 场拼成复联络 $\mathcal{A}=A+i\phi$」所给出的复结构，而平坦复联络的模空间只依赖 $C$ 的光滑结构（Atiyah–Bott 的辛商构造只用 $C$ 上的 $\bar\partial$ 算子）。所以 $(J,\Omega_{J})$ 这一侧确实不依赖 $C$ 的复结构；依赖复结构的是另一侧 $(I,\Omega_{I})$，即 Higgs 丛的复结构。}
\expandafter\def\csname zhnsrc@2.6\endcsname{[来源：A+B；nLab 词条「flat connection」（2026-11-15 访问）；Atiyah–Bott (1982)]}

\expandafter\def\csname zhnlab@2.7\endcsname{2.7}
\expandafter\def\csname zhnbody@2.7\endcsname{「一般」指避开那些使 $\MH$ 出现额外奇点、或使 Hitchin 纤维化的纤维退化的 $C$（例如分歧点重合、或 $\MH$ 上出现额外对称的情形）。在一般情形下 $\MH$ 光滑、纤维是阿贝尔簇，下面关于两个模空间不相等的结论才逐字成立。}
\expandafter\def\csname zhnsrc@2.7\endcsname{[来源：B；nLab 词条「Hitchin fibration」（2026-11-15 访问）]}

\expandafter\def\csname zhnlab@2.8\endcsname{2.8}
\expandafter\def\csname zhnbody@2.8\endcsname{$\MH(G,C)$ 与 $\MH(\LG,C)$ 的复维数相同（都是 $2(g-1)\dim G$），但并不典范同构，一般甚至不同胚。所以 B-模型与 A-模型之间的镜像对称不能被解释为某个空间上的自同构或某个变换；它是两台不同空间上的理论之间的等价，这正是第 1 节所说「不用通量也可以有镜对称」的具体体现。}
\expandafter\def\csname zhnsrc@2.8\endcsname{[来源：A+B；nLab 词条「Hitchin fibration」（2026-11-15 访问）；Hausel–Thaddeus (2002)]}

\expandafter\def\csname zhnlab@3.1\endcsname{3.1}
\expandafter\def\csname zhnbody@3.1\endcsname{$G=\SU(2)$ 时 $\phi\in H^{0}(C,K\otimes\ad E)$，而 $\ad E$ 是无迹部分，故 $\det\phi\in H^{0}(C,K^{2})$，即一个二次微分。特征多项式只有 $\det$ 一个系数（$\operatorname{tr}\phi=0$），所以 Hitchin 映射就是 $\phi\mapsto\det\phi$。$B=H^{0}(C,K^{2})$ 的维数是 $3g-3$，恰为 $\MH$ 复维数的一半，与「纤维是环面」相符。}
\expandafter\def\csname zhnsrc@3.1\endcsname{[来源：A+B；nLab 词条「spectral curve」（Last revised October 20, 2017；2026-11-15 访问）；Hitchin (1987b)]}

\expandafter\def\csname zhnlab@3.2\endcsname{3.2}
\expandafter\def\csname zhnbody@3.2\endcsname{一般地，$B=\bigoplus_{i=1}^{r}H^{0}(C,K^{d_{i}})$，$r=\operatorname{rank}G$，$d_{i}$ 是 $G$ 各基本不变多项式的次数（等于指数加一；例如 $A_{n-1}$ 时为 $2,3,\dots,n$）。这是把 $\phi$ 的特征多项式的各系数当作 $B$ 的坐标：$\operatorname{tr}\phi^{k}$ 落在 $H^{0}(C,K^{k})$ 中。$\dim B=\sum_{i}(2d_{i}-1)(g-1)=(g-1)\dim G$，正好是 $\MH$ 维数的一半；$B$ 本身是仿射空间。}
\expandafter\def\csname zhnsrc@3.2\endcsname{[来源：A+B；nLab 词条「Hitchin fibration」（2026-11-15 访问）；Hitchin (1987b)]}

\expandafter\def\csname zhnlab@3.3\endcsname{3.3}
\expandafter\def\csname zhnbody@3.3\endcsname{$G=\SU(2)$ 时一般纤维是 Prym 簇（或 Jacobian 的 torsor），一般 $G$ 时是 Hitchin 所构造的「Prym 的推广」。之所以是 torsor 而非簇，是因为纤维上没有自然的原点，除非额外取定一个截面。纤维是环面这一点，使 $\MH\to B$ 成为完全可积 Hamilton 系统：$B$ 是作用量变量，纤维是角变量。}
\expandafter\def\csname zhnsrc@3.3\endcsname{[来源：A+B；nLab 词条「integrable system」（Last revised April 29, 2026；2026-11-15 访问）；Hitchin (1987b)；Hausel–Thaddeus (2002)]}

\expandafter\def\csname zhnlab@3.4\endcsname{3.4}
\expandafter\def\csname zhnbody@3.4\endcsname{特殊拉格朗日子流形 $L\subset(Y,\Omega,\omega)$ 指 $\omega|_{L}=0$（拉格朗日）且 $\operatorname{Im}(e^{i\vartheta}\Omega)|_{L}=0$（相位固定）。这里 $L$ 关于 $I$ 全纯，故 $\Omega_{I}|_{L}$ 是 $L$ 上的全纯体积形式；又 $L$ 关于 $\Omega_{I}=\omega_{J}+i\omega_{K}$ 拉格朗日，特别地 $\omega_{K}|_{L}=0$，于是 $L$ 关于 $\omega_{K}$ 是特殊拉格朗日。复子流形在 Kähler 流形中面积最小（Wirtinger 定理），所以这些纤维在同调类中极小——原文的「面积最小」正是此意。}
\expandafter\def\csname zhnsrc@3.4\endcsname{[来源：B；nLab 词条「special Lagrangian submanifold」（2026-11-15 访问）；Hitchin (1987b)；Strominger–Yau–Zaslow (1996)]}

\expandafter\def\csname zhnlab@3.5\endcsname{3.5}
\expandafter\def\csname zhnbody@3.5\endcsname{SYZ 猜想：$X$ 与它的镜像 $X'$ 共享同一个底 $B$，两者的特殊拉格朗日环面纤维逐点互为对偶环面（T-对偶作用在纤维上）。这里把环面换成阿贝尔簇：对偶阿贝尔簇即 Pontryagin 对偶 $\operatorname{Hom}(A,\mathbb{R}/\mathbb{Z})$ 的复化。Hitchin 纤维化正好提供 SYZ 结构，而「对偶群的 Hitchin 纤维化」正是那个对偶纤维化——这就是为什么 $\MH(G,C)$ 与 $\MH(\LG,C)$ 会有同一个底 $B$。}
\expandafter\def\csname zhnsrc@3.5\endcsname{[来源：A+B；nLab 词条「T-duality」（Last revised August 14, 2026；2026-11-15 访问）；nLab 词条「mirror symmetry」（2026-11-15 访问）；Strominger–Yau–Zaslow (1996)；Harvey–Moore–Strominger (1995)]}

\expandafter\def\csname zhnlab@3.6\endcsname{3.6}
\expandafter\def\csname zhnbody@3.6\endcsname{依赖 $C$ 的复结构的是「工具」（Hitchin 方程、纤维化），不是「结论」（镜像对本身）：第 5、6 节会说明，这一镜像对称可从四维 $\mathcal{N}=4$ 理论经拓扑扭转得到，而扭转后的理论只依赖 $C$ 的定向光滑结构。换句话说，$\MH\to B$ 是证明用的脚手架，拆掉之后剩下的等价关系只认 $C$ 的二维拓扑。}
\expandafter\def\csname zhnsrc@3.6\endcsname{[来源：A+B；nLab 词条「topological twist」（Last revised October 10, 2025；2026-11-15 访问）；Kapustin–Witten (2007)]}

\expandafter\def\csname zhnlab@3.7\endcsname{3.7}
\expandafter\def\csname zhnbody@3.7\endcsname{D-模按定义是微分算子层 $\mathcal{D}$ 上的（准）凝聚模；在 $M(G,C)$ 上，$\mathcal{D}$-模就是「带平坦联络的向量丛」的无穷小版本（Riemann–Hilbert 对应）。它与 A-膜的对应是 Kapustin–Witten 第 11 节的核心：$\MH(G,C)$ 上存在一个特殊的 A-膜（共各向同性膜），其自同态代数正是 $M(G,C)$ 上微分算子的整体截面；而 $\MH(G,C)\supset T^{*}M_{\mathrm{st}}(G,C)$ 是 Zariski 开集，这个开集上的 Higgs 丛正是稳定丛的余切丛。}
\expandafter\def\csname zhnsrc@3.7\endcsname{[来源：A+B；nLab 词条「D-module」（Last revised July 27, 2024；2026-11-15 访问）；nLab 词条「Kapustin-Witten TQFT」（Last revised September 28, 2026；2026-11-15 访问）；Kapustin–Witten (2007) 第 11 节；Hitchin (1987a)]}

\expandafter\def\csname zhnlab@3.8\endcsname{3.8}
\expandafter\def\csname zhnbody@3.8\endcsname{几何 Hecke 算子在点 $x\in C$ 处对 $G$-丛作「基本修改」：把 $E$ 在 $x$ 附近平凡化、乘以 $z^{\lambda}$ 再粘回去，得到同构类可能改变的 $E'$；它诱导 $M(G,C)$ 上 D-模之间的对应，其特征层即 Hecke 特征层。Bogomolny 方程是三维规范场论中的单极子方程 $F=\star\,d_{A}\phi$，正是 Hitchin 方程在 $C=\mathbb{R}^{2}$、并沿第三方向平移不变时的约化；'t Hooft 线算子则是沿一条线的奇点条件。Kapustin–Witten 说明：Wilson 与 't Hooft 线算子在四维理论中的对偶，正好翻译成 Hecke 算子与 Hecke 特征层的陈述。}
\expandafter\def\csname zhnsrc@3.8\endcsname{[来源：A+B；nLab 词条「Hecke operator」（Last revised February 17, 2025；2026-11-15 访问）；nLab 词条「Bogomolny equation」（2026-11-15 访问）；Kapustin–Witten (2007) 第 9、10 节；Atiyah–Hitchin (1988)]}

\expandafter\def\csname zhnlab@4.1\endcsname{4.1}
\expandafter\def\csname zhnbody@4.1\endcsname{挖去一点后 $C'$ 非紧，$\pi_{1}(C')$ 是秩 $2g+1$ 的自由群（$2g$ 个生成元加一个绕 $p$ 的生成元），故 $C'$ 上的平坦丛可以有非平凡单值性；这正是「分歧」的几何来源。取 $C'$ 而不是 $C$，是因为平坦丛在 $p$ 处一般不能延拓。}
\expandafter\def\csname zhnsrc@4.1\endcsname{[来源：A+B；nLab 词条「flat connection」（2026-11-15 访问）；Simpson (1990)]}

\expandafter\def\csname zhnlab@4.2\endcsname{4.2}
\expandafter\def\csname zhnbody@4.2\endcsname{绕 $p$ 的单值性是一个共轭类：改换基点只把它共轭一下，故「单值性落在 $\mathcal{C}$ 中」是良定义的。固定 $\mathcal{C}$ 是为了让模空间有预期维数——否则单值性的类型会在模空间里跳变，得到的是分层的空间而非流形。$\mathcal{C}$ 由半单部分（极大环面中的特征值，模 Weyl 群）与幂零部分共同刻画。}
\expandafter\def\csname zhnsrc@4.2\endcsname{[来源：A+B；nLab 词条「moduli stack of G-principal bundles」（Last revised February 24, 2026；2026-11-15 访问）；Simpson (1990)]}

\expandafter\def\csname zhnlab@4.3\endcsname{4.3}
\expandafter\def\csname zhnbody@4.3\endcsname{两个理由。其一，$G_{\C}$ 的共轭类由极大环面中的特征值刻画，而 ${}^{L}G_{\C}$ 的极大环面是对偶环面（权格与余权格互换），二者之间没有典范映射把「乘性」的特征值送成特征值——只有取对数之后才有线性对偶。其二（更根本），B-模型全纯地依赖 $\mathcal{C}$（共轭类是复解析数据），而 A-模型只依赖 $\omega=\operatorname{Im}\Omega$，对 $\mathcal{C}$ 的依赖是实的；两者依赖方式不同，故不能直接相等。要得到正确的猜想必须扩充参数，即第 4 节末的 $(\alpha,\beta,\gamma,\eta)$。}
\expandafter\def\csname zhnsrc@4.3\endcsname{[来源：B；nLab 词条「Langlands dual group」（2026-11-15 访问）；nLab 词条「function field analogy」（2026-11-15 访问）；Gukov–Witten (2006) 第 2 节]}

\expandafter\def\csname zhnlab@4.4\endcsname{4.4}
\expandafter\def\csname zhnbody@4.4\endcsname{正则指中心化子维数取最小值（等于 $\operatorname{rank}G$）；对半单元，正则等价于特征值两两不同，也等价于中心化子是极大环面。原文取 $U$ 正则，是为了让 (4.1)(4.2) 中的 $\alpha,\beta,\gamma$ 干净地落在 $\mathfrak{t}$ 里、且 $U$ 的共轭类有固定维数。非正则（例如 $U=1$，即无分歧的极限）会使模空间出现额外奇点与分层，但镜像对称的陈述仍成立，只是表述更繁；一般情形的完整陈述见 Gukov 与 Witten。}
\expandafter\def\csname zhnsrc@4.4\endcsname{[来源：A；Gukov–Witten (2006)]}

\expandafter\def\csname zhnlab@4.5\endcsname{4.5}
\expandafter\def\csname zhnbody@4.5\endcsname{(4.1)(4.2) 是 Hitchin 方程在 $z=re^{i\theta}$ 极坐标下的模型解：$A$ 的角向分量给出绕 $p$ 的 holonomy 的对数，$\phi$ 的 $dr/r$ 与 $d\theta$ 分量给出 Higgs 场在穿刺点附近的极点结构。把复联络 $\mathcal{A}=A+i\phi$ 沿绕 $p$ 的小圆积分得 $\oint\mathcal{A}=2\pi(\alpha-i\gamma)$，故 holonomy 为 $\exp(2\pi(\alpha-i\gamma))$；原文写作 $U=\exp(-2\pi(\alpha-i\gamma))$，相差的符号来自绕行方向与 holonomy 取法的约定，两种写法在文献中都常见，只需在 $G\leftrightarrow{}^{L}G$ 的映射中一致取定。$\alpha$ 是「电」参数（联络），$\gamma$ 是「磁」参数（Higgs 场），$\beta$ 控制 $\phi$ 在径向上的极点强度。}
\expandafter\def\csname zhnsrc@4.5\endcsname{[来源：A+B；nLab 词条「Hitchin system」（2026-11-15 访问）；Simpson (1990)；Gukov–Witten (2006) 第 2 节]}

\expandafter\def\csname zhnlab@4.6\endcsname{4.6}
\expandafter\def\csname zhnbody@4.6\endcsname{$\eta$ 在四维 $\mathcal{N}=4$ 理论中是 $\theta$-角与耦合常数的组合（更准确地说，扭转后它成为复化 Kähler 类的 $B$-场部分）。四维的 S-对偶把 $\tau=\theta/2\pi+4\pi i/g^{2}$ 换成 $-1/\tau$，作用在 $(\alpha,\eta)$ 上正是旋转 $90^{\circ}$：$\alpha\mapsto\eta$、$\eta\mapsto-\alpha$（原文写作 $(\alpha,\eta)=({}^{L}\eta,-{}^{L}\alpha)$）。其含义是：一侧 A-模型的 Kähler 参数 $\eta$，在对偶侧变成 B-模型的复结构参数 ${}^{L}\alpha$——经典几何看不见 $\eta$，只有量子层面才看得见，这正是原文强调「单值性的对偶涉及量子参数」的原因。}
\expandafter\def\csname zhnsrc@4.6\endcsname{[来源：A+B；nLab 词条「S-duality」（Last revised September 18, 2026；2026-11-15 访问）；Gukov–Witten (2006) 第 2.3、2.4 节]}

\expandafter\def\csname zhnlab@4.7\endcsname{4.7}
\expandafter\def\csname zhnbody@4.7\endcsname{正则奇点指 $A=dz\,(M/z+\text{正则项})$、$M$ 为常值矩阵的情形；等价地，联络在 $p$ 处的奇性至多是一阶极点。驯顺分歧在代数几何中正是指这种情形（可能还要对 $C$ 取有限分歧覆盖以消去留数的非半单部分）。「驯顺分歧」一词在文献中有两种用法：一指单极点（不论留数），二指单极点且留数幂零；原文此处点明这两种用法的差别，并说明 Gukov–Witten 的框架统一处理任意留数：幂零留数对应 ${}^{L}\alpha={}^{L}\gamma=0$，半单留数对应本节 (4.3) 的情形，两者之间可以连续变形。}
\expandafter\def\csname zhnsrc@4.7\endcsname{[来源：A+B；nLab 词条「Stokes phenomenon」（Last revised October 15, 2024；2026-11-15 访问）；nLab 词条「D-module」（Last revised July 27, 2024；2026-11-15 访问）；Gukov–Witten (2006)；Simpson (1990)]}

\expandafter\def\csname zhnlab@5.1\endcsname{5.1}
\expandafter\def\csname zhnbody@5.1\endcsname{留数是 $A$ 在 $z=0$ 处 $dz/z$ 的系数。在 $z\mapsto\lambda z$ 下 $dz/z$ 不变，故 $T_{1}$ 与坐标选择无关；它是绕 $p$ 的 holonomy 的一阶（对数）信息，因此在纯拓扑意义下有意义。}
\expandafter\def\csname zhnsrc@5.1\endcsname{[来源：A+B；nLab 词条「Stokes phenomenon」（Last revised October 15, 2024；2026-11-15 访问）；Simpson (1990)]}

\expandafter\def\csname zhnlab@5.2\endcsname{5.2}
\expandafter\def\csname zhnbody@5.2\endcsname{在 $z\mapsto\lambda z$ 下，$A$ 作为 1-形式不变，而 $dz(T_{k}/z^{k})$ 变成 $dz(\lambda^{1-k}T_{k}/z^{k})$，故 $T_{k}$ 按 $\lambda^{k-1}$ 变换：$T_{1}$ 不变、$T_{2}$ 线性、$T_{n}$ 按 $\lambda^{n-1}$ 变换。所以 $T_{2},\dots,T_{n}$ 不是坐标无关的对象——单个 $T_{k}$ 没有拓扑意义，只有「$T_{n}$ 作为 $K_{p}^{n-1}$-值对象」才有（见译注 5.6、5.7）。这正是原文追问「它们怎能进入一个只依赖定向二维流形的理论」的技术根源。}
\expandafter\def\csname zhnsrc@5.2\endcsname{[译者按（由原文式 (5.1) 直接计算）]}

\expandafter\def\csname zhnlab@5.3\endcsname{5.3}
\expandafter\def\csname zhnbody@5.3\endcsname{$d\Phi+\Phi^{2}=0$ 是把 $(A,\phi)$ 拼成 $\Phi$ 之后「曲率为零」的简写（真正的 Hitchin 方程还含 $D\star\phi=0$）。量级比较是：$\Phi\sim|z|^{-n}$ 时 $d\Phi\sim|z|^{-(n+1)}$、$\Phi^{2}\sim|z|^{-2n}$；$n=1$ 时二者同阶，方程平衡，可用隐函数定理在 $0$ 附近求解；$n>1$ 时 $\Phi^{2}$ 严格主导，线性化方程的解无法微扰地延拓到 $z=0$，故说「非线性项太强」。这不是说方程无解，而是说不能用通常的隐函数或逐次逼近法——出路是本节后半段的交换化。}
\expandafter\def\csname zhnsrc@5.3\endcsname{[来源：A+B；nLab 词条「Hitchin system」（2026-11-15 访问）；Hitchin (1987a)；Biquard–Boalch (2004)]}

\expandafter\def\csname zhnlab@5.4\endcsname{5.4}
\expandafter\def\csname zhnbody@5.4\endcsname{非正则奇点的形式解只在若干条 Stokes 射线上有渐近展开，跨过射线时展开式要乘上一个常数矩阵，即 Stokes 矩阵。Stokes 矩阵连同形式单值性合称扩张单值性；它只依赖 $C$ 的定向、奇点阶数与留数的共轭类，因此是纯拓扑数据——这正是「高阶极点也有拓扑含义」的答案。}
\expandafter\def\csname zhnsrc@5.4\endcsname{[来源：A+B；nLab 词条「Stokes phenomenon」（Last revised October 15, 2024；2026-11-15 访问）；Wasow (1965)；Boalch (2001)]}

\expandafter\def\csname zhnlab@5.5\endcsname{5.5}
\expandafter\def\csname zhnbody@5.5\endcsname{等单值形变指保持（扩张）单值数据不变、只改变 $T_{2},\dots,T_{n}$ 的形变；其形变方程是 Painlevé 型方程，Jimbo–Miwa–Ueno 证明它们是 Hamilton 系统。Boalch 进一步给出这套 Hamilton 系统的辛结构，并证明它沿等单值形变流不变；于是可以把 $T_{2},\dots,T_{n}$ 方向商掉，得到只依赖 $T_{1}$ 与 $n$ 的复辛流形 $\YC(G,C;T_{1},n)$。注意此处只得到复辛结构，还没有度量或超凯勒结构——后者要靠下一段的 Hitchin 方程。}
\expandafter\def\csname zhnsrc@5.5\endcsname{[来源：A+B；nLab 词条「isomonodromic deformation」（Last revised January 13, 2021；2026-11-15 访问）；Jimbo–Miwa–Ueno (1981)；Boalch (2001)]}

\expandafter\def\csname zhnlab@5.6\endcsname{5.6}
\expandafter\def\csname zhnbody@5.6\endcsname{$\mathfrak{t}_{\C}^{\mathrm{reg}}$ 是 $\mathfrak{t}_{\C}$ 中正则元（特征值两两不同）构成的开集；$\mathfrak{t}_{\C}^{\mathrm{reg}}\otimes K_{p}^{n-1}$ 表示取值于 $\mathfrak{t}_{\C}^{\mathrm{reg}}$ 的 $K_{p}^{n-1}$-值对象，它把译注 5.2 的缩放律 $\lambda^{n-1}$ 吸收进张量因子，使 $T_{n}$ 成为坐标无关的对象。辫群：原文说 $\mathfrak{t}_{\C}^{\mathrm{reg}}$ 的基本群，严格说这是纯辫群；文献中「$G$ 的辫群」通常指 $\pi_{1}(\mathfrak{t}_{\C}^{\mathrm{reg}}/W)$（$W$ 为 Weyl 群），二者相差一个 $W$。原文的用法（以及 $B(G)\cong B({}^{L}G)$ 的陈述）对这两种取法都成立，此处按原文照译。}
\expandafter\def\csname zhnsrc@5.6\endcsname{[来源：A+B；nLab 词条「braid group」（Last revised August 17, 2026；2026-11-15 访问）；Boalch (2001)]}

\expandafter\def\csname zhnlab@5.7\endcsname{5.7}
\expandafter\def\csname zhnbody@5.7\endcsname{因为 $T_{n}$ 取值于 $\mathfrak{t}_{\C}^{\mathrm{reg}}\otimes K_{p}^{n-1}$，而 $K_{p}$ 是复线、绕 $p$ 的 Dehn twist（映射类群的一部分）作用在 $K_{p}$ 上带一个相位；这个相位又与 $T_{n}$ 沿 $\mathfrak{t}_{\C}^{\mathrm{reg}}$ 绕行时的辫群单值性耦合，故整体的对称群不是「映射类群 $\times\,B(G)$」的直积，而是映射类群被 $B(G)$ 的扩张（半直积或中心扩张）。原文的 extension of the mapping class group by $B(G)$ 即指此。}
\expandafter\def\csname zhnsrc@5.7\endcsname{[来源：A+B；nLab 词条「mapping class group」（Last revised August 11, 2026；2026-11-15 访问）；nLab 词条「braid group」（Last revised August 17, 2026；2026-11-15 访问）；Boalch (2001)]}

\expandafter\def\csname zhnlab@6.1\endcsname{6.1}
\expandafter\def\csname zhnbody@6.1\endcsname{规范场是 1-形式联络 $A$，其曲率 $F$ 是 2-形式；电磁对偶把 $F$ 换成 Hodge 对偶 $\star F$。在 $d$ 维中 $\star$ 把 $p$-形式送到 $(d-p)$-形式，故 2-形式的对偶仍是 2-形式当且仅当 $d=4$。$d\neq4$ 时对偶把规范场送到 $(d-3)$-形式场（「磁」规范场是高阶形式规范场），理论类型随之改变——这就是为什么规范场论的对偶最自然地在四维发生。}
\expandafter\def\csname zhnsrc@6.1\endcsname{[来源：B；nLab 词条「S-duality」（Last revised September 18, 2026；2026-11-15 访问）；nLab 词条「electro-magnetic duality」（Last revised July 24, 2026；2026-11-15 访问）]}

\expandafter\def\csname zhnlab@6.2\endcsname{6.2}
\expandafter\def\csname zhnbody@6.2\endcsname{$\mathcal{N}=4$ 超 Yang–Mills 在四维有 16 个超荷，是无引力情形下可能的最大超对称。Montonen 与 Olive 猜想它有 S-对偶：耦合 $g\leftrightarrow1/g$、规范群 $G\leftrightarrow{}^{L}G$，磁单极子与规范粒子互换。拓扑扭转指把 Lorentz 群换成「Lorentz 与 R-对称的对角子群」，使某个超荷变成标量，从而可以在一类流形上定义拓扑不变量；$\mathcal{N}=2$ 的相应扭转给出 Donaldson 理论，$\mathcal{N}=4$ 的扭转则给出 Kapustin–Witten 的两族拓扑场论，正好对应 A-模型与 B-模型。}
\expandafter\def\csname zhnsrc@6.2\endcsname{[来源：A+B；nLab 词条「topological twist」（Last revised October 10, 2025；2026-11-15 访问）；nLab 词条「S-duality」（Last revised September 18, 2026；2026-11-15 访问）；Brink–Schwarz–Scherk (1977)；Montonen–Olive (1977)；Kapustin–Witten (2007)]}

\expandafter\def\csname zhnlab@6.3\endcsname{6.3}
\expandafter\def\csname zhnbody@6.3\endcsname{$\sigma$ 模型是「从世界面 $\Sigma$ 到目标 $X$ 的映射」的量子场论，作用量由 $\Sigma$ 上的度量拉回与 $B$-场组成。$\mathcal{N}=(2,2)$ 超对称的 $\sigma$ 模型有两种拓扑扭转：A-扭转得到 A-模型（关联函数计数伪全纯曲线），B-扭转得到 B-模型（关联函数是复结构变形的积分）。$\Sigma$ 是辅助的：理论的物理内容不依赖 $\Sigma$ 的度量与复结构；拓扑扭转之后连 $\Sigma$ 的复结构也不再进入。}
\expandafter\def\csname zhnsrc@6.3\endcsname{[来源：B；nLab 词条「A-model」（Last revised October 23, 2026；2026-11-15 访问）；nLab 词条「B-model」（Last revised February 12, 2025；2026-11-15 访问）；nLab 词条「topological twist」（Last revised October 10, 2025；2026-11-15 访问）；Hori 等 (2003)]}

\expandafter\def\csname zhnlab@6.4\endcsname{6.4}
\expandafter\def\csname zhnbody@6.4\endcsname{提升的要点：把 $\MH(G,C)$ 上的点 $(A(y),\phi(y))$ 换成同时依赖 $x\in\Sigma$ 的场 $(A(x,y),\phi(x,y))$，即允许 Higgs 丛沿 $\Sigma$ 方向变化。此时 $A$ 不再只是「切于 $C$ 的 1-形式」，而可以是 $M=\Sigma\times C$ 上一般的 1-形式；$\phi$ 也从 $C$ 上的 $(1,0)$-形式延拓为 $M$ 上的 $\ad$-值 1-形式。把这组场按最小方式补成超对称多重态，得到的正是四维 $\mathcal{N}=4$ 理论在 $M$ 上的扭转版本——它比 $\sigma$ 模型多了沿 $\Sigma$ 方向的动能项与超对称伙伴。}
\expandafter\def\csname zhnsrc@6.4\endcsname{[来源：A+B；nLab 词条「topological twist」（Last revised October 10, 2025；2026-11-15 访问）；Kapustin–Witten (2007) 第 3 节]}

\expandafter\def\csname zhnlab@6.5\endcsname{6.5}
\expandafter\def\csname zhnbody@6.5\endcsname{$\sigma$ 模型只在 $\MH(G,C)$ 的光滑点上有定义；遇到奇点，映射 $\Phi$ 的模空间本身退化，「计数」无从谈起。四维理论则在 $M=\Sigma\times C$ 上处处良定义，因为它不把 $\MH(G,C)$ 当作目标空间。两者的关系是「零模约化」：当 $\Sigma$ 的尺度远大于 $C$ 时，四维理论沿 $\Sigma$ 的涨落被压低，剩下的正是 $\Sigma\to\MH(G,C)$ 的 $\sigma$ 模型；在奇点附近这一约化失效。}
\expandafter\def\csname zhnsrc@6.5\endcsname{[来源：A+B；nLab 词条「topological twist」（Last revised October 10, 2025；2026-11-15 访问）；Kapustin–Witten (2007)]}

\expandafter\def\csname zhnlab@6.6\endcsname{6.6}
\expandafter\def\csname zhnbody@6.6\endcsname{Hecke 算子在点 $x$ 处修改丛，会改变丛的斜率（次数与秩的比值），因此会把稳定丛送到不稳定丛；若只在半稳定模空间 $M(G,C)$ 上取 D-模，Hecke 算子就不保持这个空间，无法定义其作用。所有 $G$-丛的叠 $\Bun_{G}(C)$ 才是 Hecke 算子能自由作用的地方；$M(G,C)$ 只是其中「稳定」的开子集。不稳定丛正对应 $M(G,C)$ 的非轨形奇点。}
\expandafter\def\csname zhnsrc@6.6\endcsname{[来源：A+B；nLab 词条「Hecke operator」（Last revised February 17, 2025；2026-11-15 访问）；nLab 词条「Hecke eigensheaf」（2026-11-15 访问）；Kapustin–Witten (2007) 第 9 节]}

\expandafter\def\csname zhnlab@6.7\endcsname{6.7}
\expandafter\def\csname zhnbody@6.7\endcsname{叠可以理解为「局部看起来像商」的空间（用 groupoid 值函子或局部商来严格化）。有限维商不存在的理由是一个维数计数：$G_{\C}$-丛 $E$ 越不稳定，其自同构群越大（例如 $E$ 含大量平凡直和项时，自同构含任意多个多项式截面），自同构群的维数没有上界；而有限维 $W_{\C}$ 的稳定子维数有上界。故只能取无穷维的 $(Y,W_{\C})$。}
\expandafter\def\csname zhnsrc@6.7\endcsname{[来源：A+B；nLab 词条「moduli stack」（Last revised November 14, 2026；2026-11-15 访问）；Atiyah–Bott (1982)]}

\expandafter\def\csname zhnlab@6.8\endcsname{6.8}
\expandafter\def\csname zhnbody@6.8\endcsname{作用量中的 $|\mu|^{2}$ 项（$\mu$ 为 $W$ 作用的矩映射）在 $\mu=0$ 时取极小；$\mu^{-1}(0)/W$ 与 GIT 商 $Y/\!/W_{\C}$ 在经典层面一致，这就是 Kempf–Ness 定理的内容：GIT 商等于辛商。物理上这是「D-项」与「F-项」的关系；量子层面还要加上 $\theta$-角，故原文说「经典上等价」。}
\expandafter\def\csname zhnsrc@6.8\endcsname{[来源：B；nLab 词条「symplectic quotient」（Last revised September 20, 2013；2026-11-15 访问）；Atiyah–Bott (1982)]}

\expandafter\def\csname zhnlab@6.9\endcsname{6.9}
\expandafter\def\csname zhnbody@6.9\endcsname{不变正定二次型存在等价于群紧：若 $W$ 紧，可对 Haar 测度平均任意正定型得到不变型；反之若存在非紧方向，则不存在这样的型。无穷维的 $W=\Maps(C,G)$ 之所以仍有这样的型，是因为 $G$ 紧、$\mathfrak{g}$ 上有 $\operatorname{Ad}$-不变正定型，再用 $C$ 上的积分（需要 $C$ 上的测度，例如黎曼面的面积测度）把它逐点积分起来。}
\expandafter\def\csname zhnsrc@6.9\endcsname{[来源：B；nLab 词条「symplectic quotient」（Last revised September 20, 2013；2026-11-15 访问）；nLab 词条「Langlands dual group」（2026-11-15 访问）]}

\expandafter\def\csname zhnlab@6.10\endcsname{6.10}
\expandafter\def\csname zhnbody@6.10\endcsname{四元数 Lie 群（李代数带三个满足四元数关系、且与括号相容的复结构）在非平凡情形下不存在，所以不能像辛商那样用一个四元数群取商。超凯勒商绕开这一点：它同时取三个矩映射 $\mu_{I},\mu_{J},\mu_{K}$，以 $\mu^{-1}(0)/W$ 代替辛商，得到的仍是超凯勒流形。Hitchin 正是用它把 $\MH(G,C)$ 实现为「$C$ 上 $(A,\phi)$ 对」的空间被 $\Maps(C,G)$ 的超凯勒商。代数几何意义上的 Higgs 丛叠只记住复结构 $I$（从而只记住 $\Omega_{I}$ 与 $\mu_{I}$），丢掉了 $J,K$ 与另外两个矩映射，故原文说它「过于强调一个复结构」、不足以表达这里的超凯勒商。}
\expandafter\def\csname zhnsrc@6.10\endcsname{[来源：A+B；nLab 词条「hyperkähler manifold」（Last revised July 14, 2021；2026-11-15 访问）；nLab 词条「symplectic quotient」（Last revised September 20, 2013；2026-11-15 访问）；Hitchin (1987a)；Hitchin–Karlhede–Lindström–Roček (1987)]}
